\documentclass[11pt,paper=a4,answers]{exam}
\usepackage{graphicx,lastpage}
\usepackage{upgreek}
\usepackage{censor}
\usepackage{gensymb}
\usepackage{amsmath}
\usepackage{multicol}
\usepackage{enumitem}
\usepackage{tasks}
\usepackage{adjustbox}
\usepackage{xcolor}
\usepackage{tfrupee}
\setlength{\voffset}{0.25in}
\usepackage{tabularx}
\newcolumntype{Y}{>{\centering\arraybackslash}X}
%==============================================================
% Customise Non-Essential Packages Below
% =============================================================
\usepackage[most]{tcolorbox}
\usepackage{eso-pic}
\usepackage{tikz}
\AddToShipoutPictureBG{%
\begin{tikzpicture}[remember picture,overlay]
\foreach \x in {-8,-4,0,4,8,12,16,20}
\foreach \y in {-8,-4,0,4,8,12,16,20} {
\node[opacity=0.05,rotate=45,anchor=center] at ([xshift=\x cm,yshift=\y cm]current page.center) {%
\begin{minipage}{2cm}
\centering
\includegraphics[width=2cm]{images/logo_black.png}\\
\small preppeo.com
\end{minipage}%
};
}
\end{tikzpicture}%
}
\printanswers

%==============================================================
% Customise Non-Essential Packages Above
% =============================================================
\definecolor{svgray}{rgb}{0.90, 0.90, 0.90}
\graphicspath{{images/}}

\usepackage[normalem]{ulem}
\renewcommand\ULthickness{1pt}   %%---> For changing thickness of underline
\setlength\ULdepth{3.5ex}%\maxdimen ---> For changing depth of underline
\renewcommand{\baselinestretch}{1}
\chead{
\includegraphics[scale=0.1,height=2cm ]{logo.png}
}
\pagestyle{empty}

\pagestyle{headandfoot}
\headrule
\cfoot{W: https://preppeo.com; M: +91-8130711689}

\begin{document}
%==============================================================
% Customise Question paper details below this
% =============================================================
\begin{center}
    \centering
    \renewcommand{\arraystretch}{1.5}
    \begin{tabularx}{\textwidth}{YYY}
    & \normalsize\bf{{\Large{{Quantitative Reasoning}}}} & \\
    By Preppeo & \normalsize\bf{{sat\_lid\_014}} & SAT\\
    \normalsize{{\textbf{{Unit:}} Advanced Math}} & \normalsize{{\textbf{{Chapter:}} Nonlinear equations in one variable}} & \normalsize{{\textbf{{Topic:}} Solving Quadratics (Factoring \& Formula)}} \\
    \end{tabularx}
\end{center}
\par\noindent\rule{\textwidth}{1pt}
% =============================================================
% Customise Question paper details above this
% =============================================================

\begin{questions}
\pointsinrightmargin
\pointsdroppedatright
\marksnotpoints
\pointpoints{mark}{marks}
\pointformat{\boldmath\themarginpoints}
\bracketedpoints

% =============================================================
% Question Begin
% =============================================================
% Lid Topic: Advanced Math — Nonlinear Equations: Solving Quadratics
% Total Questions: 50

% --- PART 1: Factoring and Fundamental Solving (1-25) ---

% Question 1
% Category: Practice | Difficulty: Medium | Question Type: Multiple Choice
\question
Which of the following are the solutions to the equation $x^2 - 7x + 10 = 0$?
\begin{tasks}(4)
    \task $x = -2$ and $x = -5$
    \task $x = 2$ and $x = 5$
    \task $x = -10$ and $x = 1$
    \task $x = 10$ and $x = -1$
\end{tasks}

\begin{solution}
\textbf{Conceptual Explanation:} \\
To solve a quadratic equation in the form $x^2 + bx + c = 0$ by factoring, we look for two numbers that multiply to $c$ (the constant) and add to $b$ (the coefficient of $x$).

\vspace{15pt}

\textbf{Calculation and Logic:} \\
1. Identify the factors of 10 that add up to $-7$. \\
   Factors of 10: $(-1, -10)$ and $(-2, -5)$. \\
   Sum: $-2 + (-5) = -7$.

\vspace{10pt}

2. Rewrite the equation in factored form: \\
   $(x - 2)(x - 5) = 0$

\vspace{10pt}

3. Set each factor to zero using the Zero Product Property: \\
   $x - 2 = 0 \Rightarrow x = 2$ \\
   $x - 5 = 0 \Rightarrow x = 5$

\vspace{20pt}

Answer: \colorbox{yellow!100}{B}
\end{solution}

\vspace{30pt}

% Question 2
% Category: Practice | Difficulty: Hard | Question Type: Multiple Choice
\question
What is the sum of the solutions to the equation $3x^2 - 12x + 9 = 0$?
\begin{tasks}(4)
    \task 3
    \task 4
    \task 9
    \task 12
\end{tasks}

\begin{solution}
\textbf{Conceptual Explanation:} \\
While you can solve the equation fully, the SAT often tests the "Sum of Roots" theorem. For any quadratic $ax^2 + bx + c = 0$, the sum of the solutions is given by $-\frac{b}{a}$.

\vspace{15pt}

\textbf{Calculation and Logic:} \\
\textbf{Method 1: Using the Theorem} \\
$a = 3, b = -12$ \\
Sum $= -\frac{-12}{3} = \frac{12}{3} = 4$

\vspace{10pt}

\textbf{Method 2: Solving by Factoring} \\
Divide the entire equation by 3: \\
$x^2 - 4x + 3 = 0$ \\
$(x - 3)(x - 1) = 0$ \\
Solutions: $x = 3, x = 1$ \\
Sum: $3 + 1 = 4$

\vspace{20pt}

Answer: \colorbox{yellow!100}{B}
\end{solution}

\vspace{30pt}

% Question 3
% Category: Practice | Difficulty: Hard | Question Type: Multiple Choice
\question
For which of the following equations is $x = -3$ a solution?
\begin{tasks}(2)
    \task $x^2 - 9 = 0$
    \task $x^2 + 6x + 9 = 0$
    \task $x^2 - 6x + 9 = 0$
    \task $x^2 + 3 = 0$
\end{tasks}

\begin{solution}
\textbf{Conceptual Explanation:} \\
A value is a solution to an equation if substituting that value for $x$ makes the equation true ($0 = 0$).

\vspace{15pt}

\textbf{Calculation and Logic:} \\
Test Choice B: $x^2 + 6x + 9 = 0$ \\
Substitute $x = -3$: \\
$(-3)^2 + 6(-3) + 9 = 0$ \\
$9 - 18 + 9 = 0$ \\
$0 = 0$

\vspace{10pt}

Choice B is a "Perfect Square Trinomial" which factors to $(x+3)^2 = 0$, giving the root $x = -3$.

\vspace{20pt}

Answer: \colorbox{yellow!100}{B}
\end{solution}

\vspace{30pt}

% Question 4
% Category: Practice | Difficulty: Hard | Question Type: Multiple Choice
\question
If $x > 0$ and $2x^2 + 5x - 12 = 0$, what is the value of $x$?
\begin{tasks}(4)
    \task 1.5
    \task 4
    \task 3
    \task 2.5
\end{tasks}

\begin{solution}
\textbf{Conceptual Explanation:} \\
When the leading coefficient ($a$) is not 1, we can use the "ac method" or the quadratic formula. Here, we will factor by grouping.

\vspace{15pt}

\textbf{Calculation and Logic:} \\
1. Multiply $a \times c$: $2 \times (-12) = -24$. \\
   Find factors of $-24$ that add to 5: $(8, -3)$.

\vspace{10pt}

2. Split the middle term: \\
   $2x^2 + 8x - 3x - 12 = 0$

\vspace{10pt}

3. Factor by grouping: \\
   $2x(x + 4) - 3(x + 4) = 0$ \\
   $(2x - 3)(x + 4) = 0$

\vspace{10pt}

4. Solve for $x$: \\
   $2x - 3 = 0 \Rightarrow x = 1.5$ \\
   $x + 4 = 0 \Rightarrow x = -4$

\vspace{10pt}

Since the problem states $x > 0$, the only valid solution is 1.5.

\vspace{20pt}

Answer: \colorbox{yellow!100}{A}
\end{solution}

\vspace{30pt}

% Question 5
% Category: Practice | Difficulty: Medium | Question Type: Multiple Choice
\question
What are the solutions to $(x - 4)^2 = 49$?
\begin{tasks}(4)
    \task $x = 11$ and $x = -3$
    \task $x = 53$ and $x = -45$
    \task $x = 7$ and $x = -7$
    \task $x = 11$ and $x = 3$
\end{tasks}

\begin{solution}
\textbf{Conceptual Explanation:} \\
When an equation is in the form $(\text{expression})^2 = \text{constant}$, the most efficient method is to take the square root of both sides. Remember to include both the positive and negative roots.

\vspace{15pt}

\textbf{Calculation and Logic:} \\
1. Take the square root of both sides: \\
   $\sqrt{(x - 4)^2} = \pm \sqrt{49}$ \\
   $x - 4 = \pm 7$

\vspace{10pt}

2. Solve the two resulting linear equations: \\
   Case 1: $x - 4 = 7 \Rightarrow x = 11$ \\
   Case 2: $x - 4 = -7 \Rightarrow x = -3$

\vspace{20pt}

Answer: \colorbox{yellow!100}{A}
\end{solution}

\vspace{30pt}

% Question 6
% Category: Practice | Difficulty: Hard | Question Type: Multiple Choice
\question
If the quadratic formula $x = \frac{-b \pm \sqrt{b^2 - 4ac}}{2a}$ is used to solve $x^2 - 6x + 4 = 0$, which of the following is a solution?
\begin{tasks}(2)
    \task $3 + \sqrt{5}$
    \task $6 + \sqrt{20}$
    \task $3 - \sqrt{10}$
    \task $3 + \sqrt{20}$
\end{tasks}

\begin{solution}
\textbf{Conceptual Explanation:} \\
The quadratic formula is used when a quadratic cannot be easily factored. It requires identifying $a$, $b$, and $c$ from the standard form $ax^2 + bx + c = 0$.

\vspace{15pt}

\textbf{Calculation and Logic:} \\
1. Identify coefficients: $a = 1, b = -6, c = 4$.

\vspace{10pt}

2. Substitute into the formula: \\
   $x = \frac{-(-6) \pm \sqrt{(-6)^2 - 4(1)(4)}}{2(1)}$ \\
   $x = \frac{6 \pm \sqrt{36 - 16}}{2}$ \\
   $x = \frac{6 \pm \sqrt{20}}{2}$

\vspace{10pt}

3. Simplify the radical: \\
   $\sqrt{20} = \sqrt{4 \times 5} = 2\sqrt{5}$ \\
   $x = \frac{6 \pm 2\sqrt{5}}{2}$

\vspace{10pt}

4. Divide both terms in the numerator by 2: \\
   $x = 3 \pm \sqrt{5}$

\vspace{20pt}

Answer: \colorbox{yellow!100}{A}
\end{solution}

\vspace{30pt}

% Question 7
% Category: Practice | Difficulty: Medium | Question Type: Multiple Choice
\question
How many distinct real solutions does the equation $x^2 + 10x + 25 = 0$ have?
\begin{tasks}(4)
    \task Zero
    \task One
    \task Two
    \task Infinitely many
\end{tasks}

\begin{solution}
\textbf{Conceptual Explanation:} \\
The number of real solutions is determined by the discriminant, $D = b^2 - 4ac$.
If $D > 0$, there are two solutions. 
If $D = 0$, there is one solution. 
If $D < 0$, there are no real solutions.

\vspace{15pt}

\textbf{Calculation and Logic:} \\
1. Identify coefficients: $a = 1, b = 10, c = 25$.

\vspace{10pt}

2. Calculate the discriminant: \\
   $D = (10)^2 - 4(1)(25)$ \\
   $D = 100 - 100 = 0$

\vspace{10pt}

3. Interpret the result: \\
   Since $D = 0$, there is exactly one distinct real solution (a "double root").

\vspace{20pt}

Answer: \colorbox{yellow!100}{B}
\end{solution}

% Question 8
% Category: Practice | Difficulty: Hard | Question Type: Multiple Choice
\question
Which of the following equations has the same solutions as $x^2 - 8x - 20 = 0$?
\begin{tasks}(2)
    \task $(x - 4)^2 = 36$
    \task $(x - 4)^2 = 4$
    \task $(x + 4)^2 = 36$
    \task $(x - 8)^2 = 20$
\end{tasks}

\begin{solution}
\textbf{Conceptual Explanation:} \\
This question tests "Completing the Square." The SAT often asks you to rewrite a standard quadratic equation into vertex form or a "squared binomial" form.

\vspace{15pt}

\textbf{Calculation and Logic:} \\
1. Move the constant to the right side: \\
   $x^2 - 8x = 20$

\vspace{10pt}

2. Find the value to complete the square ($(\frac{b}{2})^2$): \\
   $(\frac{-8}{2})^2 = (-4)^2 = 16$

\vspace{10pt}

3. Add 16 to both sides of the equation: \\
   $x^2 - 8x + 16 = 20 + 16$ \\
   $(x - 4)^2 = 36$

\vspace{20pt}

Answer: \colorbox{yellow!100}{A}
\end{solution}

\vspace{30pt}

% Question 9
% Category: Practice | Difficulty: Medium | Question Type: Multiple Choice
\question
If $(x - 5)(x + 3) = 0$ and $x < 0$, what is the value of $x + 5$?
\begin{tasks}(4)
    \task 2
    \task 5
    \task 8
    \task 10
\end{tasks}

\begin{solution}
\textbf{Conceptual Explanation:} \\
This is a multi-step SAT problem. First, solve the quadratic for its two possible roots, then apply the given constraint ($x < 0$), and finally perform the requested operation.

\vspace{15pt}

\textbf{Calculation and Logic:} \\
1. Solve for $x$: \\
   $x - 5 = 0 \Rightarrow x = 5$ \\
   $x + 3 = 0 \Rightarrow x = -3$

\vspace{10pt}

2. Apply the constraint $x < 0$: \\
   The only valid value for $x$ is $-3$.

\vspace{10pt}

3. Final Calculation: \\
   Substitute $x = -3$ into $x + 5$: \\
   $-3 + 5 = 2$

\vspace{20pt}

Answer: \colorbox{yellow!100}{A}
\end{solution}

% Question 10
% Category: Practice | Difficulty: Medium | Question Type: Multiple Choice
\question
What is the positive solution to the equation $2x^2 - 32 = 0$?
\begin{tasks}(4)
    \task 4
    \task 8
    \task 16
    \task 2
\end{tasks}

\begin{solution}
\textbf{Conceptual Explanation:} \\
When a quadratic equation lacks a linear term ($bx$), we can solve by isolating the $x^2$ term and taking the square root.

\vspace{15pt}

\textbf{Calculation and Logic:} \\
1. Add 32 to both sides: \\
   $2x^2 = 32$

\vspace{10pt}

2. Divide by 2: \\
   $x^2 = 16$

\vspace{10pt}

3. Take the square root: \\
   $x = \pm 4$

\vspace{10pt}

4. The question asks for the positive solution only.

\vspace{20pt}

Answer: \colorbox{yellow!100}{A}
\end{solution}

\vspace{30pt}

% Question 11
% Category: Practice | Difficulty: Hard | Question Type: Multiple Choice
\question
Which of the following is a factor of the expression $x^2 - 11x + 24$?
\begin{tasks}(4)
    \task $x + 3$
    \task $x - 6$
    \task $x - 8$
    \task $x + 8$
\end{tasks}

\begin{solution}
\textbf{Conceptual Explanation:} \\
To find the factors, we look for two integers that multiply to 24 and add to $-11$.

\vspace{15pt}

\textbf{Calculation and Logic:} \\
1. Identify factor pairs of 24: \\
   $(1, 24), (2, 12), (3, 8), (4, 6)$

\vspace{10pt}

2. Since the sum must be negative ($-11$) and the product positive ($+24$), both factors must be negative: \\
   $(-3) + (-8) = -11$

\vspace{10pt}

3. The factored form is $(x - 3)(x - 8)$. \\
   Comparing this to the choices, $x - 8$ is the correct factor.

\vspace{20pt}

Answer: \colorbox{yellow!100}{C}
\end{solution}

\vspace{30pt}

% Question 12
% Category: Practice | Difficulty: Medium | Question Type: Multiple Choice
\question
If $x^2 + kx + 16 = (x + 4)^2$, what is the value of $k$?
\begin{tasks}(4)
    \task 4
    \task 8
    \task 16
    \task 32
\end{tasks}

\begin{solution}
\textbf{Conceptual Explanation:} \\
A perfect square trinomial follows the pattern $(x + a)^2 = x^2 + 2ax + a^2$. By expanding the binomial, we can find the missing coefficient.

\vspace{15pt}

\textbf{Calculation and Logic:} \\
1. Expand the right side: \\
   $(x + 4)(x + 4) = x^2 + 4x + 4x + 16$ \\
   $= x^2 + 8x + 16$

\vspace{10pt}

2. Compare the coefficients of $x$ to the original expression $x^2 + kx + 16$: \\
   $k = 8$

\vspace{20pt}

Answer: \colorbox{yellow!100}{B}
\end{solution}

\vspace{30pt}

% Question 13
% Category: Practice | Difficulty: Hard | Question Type: Multiple Choice
\question
How many real solutions does the equation $2x^2 - 4x + 5 = 0$ have?
\begin{tasks}(4)
    \task 0
    \task 1
    \task 2
    \task Infinitely many
\end{tasks}

\begin{solution}
\textbf{Conceptual Explanation:} \\
The discriminant $b^2 - 4ac$ determines the number of real roots. If the value is negative, the roots are imaginary, meaning there are no real solutions.

\vspace{15pt}

\textbf{Calculation and Logic:} \\
1. Identify $a = 2, b = -4, c = 5$.

\vspace{10pt}

2. Substitute into the discriminant formula: \\
   $D = (-4)^2 - 4(2)(5)$ \\
   $D = 16 - 40$ \\
   $D = -24$

\vspace{10pt}

3. Since $D < 0$, the equation has zero real solutions.

\vspace{20pt}

Answer: \colorbox{yellow!100}{A}
\end{solution}

\vspace{30pt}

% Question 14
% Category: Practice | Difficulty: Medium | Question Type: Multiple Choice
\question
What is the solution set for the equation $x(x - 9) = 0$?
\begin{tasks}(4)
    \task $\{9\}$
    \task $\{0, -9\}$
    \task $\{0, 9\}$
    \task $\{-9, 9\}$
\end{tasks}

\begin{solution}
\textbf{Conceptual Explanation:} \\
The Zero Product Property states that if the product of two factors is zero, at least one of the factors must be zero.

\vspace{15pt}

\textbf{Calculation and Logic:} \\
1. Set the first factor to zero: \\
   $x = 0$

\vspace{10pt}

2. Set the second factor to zero: \\
   $x - 9 = 0 \Rightarrow x = 9$

\vspace{10pt}

The solutions are 0 and 9.

\vspace{20pt}

Answer: \colorbox{yellow!100}{C}
\end{solution}

\vspace{30pt}

% Question 15
% Category: Practice | Difficulty: Hard | Question Type: Multiple Choice
\question
The equation $x^2 + bx + 12 = 0$ has $x = 2$ as one of its solutions. What is the value of $b$?
\begin{tasks}(4)
    \task -6
    \task -8
    \task 6
    \task 8
\end{tasks}

\begin{solution}
\textbf{Conceptual Explanation:} \\
If a value is a solution to an equation, the equation must remain true when that value is substituted for $x$.

\vspace{15pt}

\textbf{Calculation and Logic:} \\
1. Substitute $x = 2$ into the equation: \\
   $(2)^2 + b(2) + 12 = 0$

\vspace{10pt}

2. Simplify and solve for $b$: \\
   $4 + 2b + 12 = 0$ \\
   $2b + 16 = 0$ \\
   $2b = -16$ \\
   $b = -8$

\vspace{20pt}

Answer: \colorbox{yellow!100}{B}
\end{solution}

\vspace{30pt}

% Question 16
% Category: Practice | Difficulty: Medium | Question Type: Multiple Choice
\question
Which of the following is equivalent to $(2x - 3)(x + 4)$?
\begin{tasks}(2)
    \task $2x^2 + 5x - 12$
    \task $2x^2 - 12$
    \task $2x^2 + 11x - 12$
    \task $2x^2 + 5x + 12$
\end{tasks}

\begin{solution}
\textbf{Conceptual Explanation:} \\
To expand two binomials, use the FOIL method (First, Outer, Inner, Last).

\vspace{15pt}

\textbf{Calculation and Logic:} \\
1. First: $2x \cdot x = 2x^2$ \\
2. Outer: $2x \cdot 4 = 8x$ \\
3. Inner: $-3 \cdot x = -3x$ \\
4. Last: $-3 \cdot 4 = -12$

\vspace{10pt}

5. Combine like terms: \\
   $2x^2 + 8x - 3x - 12$ \\
   $= 2x^2 + 5x - 12$

\vspace{20pt}

Answer: \colorbox{yellow!100}{A}
\end{solution}

\vspace{30pt}

% Question 17
% Category: Practice | Difficulty: Hard | Question Type: Multiple Choice
\question
What is the product of the roots of the equation $2x^2 - 10x + 16 = 0$?
\begin{tasks}(4)
    \task 5
    \task 8
    \task 10
    \task 16
\end{tasks}

\begin{solution}
\textbf{Conceptual Explanation:} \\
For a quadratic equation $ax^2 + bx + c = 0$, the product of the roots is given by the formula $\frac{c}{a}$.

\vspace{15pt}

\textbf{Calculation and Logic:} \\
1. Identify $a = 2$ and $c = 16$.

\vspace{10pt}

2. Apply the formula: \\
   Product $= \frac{16}{2} = 8$

\vspace{20pt}

Answer: \colorbox{yellow!100}{B}
\end{solution}

\vspace{30pt}

% Question 18
% Category: Practice | Difficulty: Medium | Question Type: Multiple Choice
\question
What are the solutions to $x^2 - 100 = 0$?
\begin{tasks}(4)
    \task $x = 10$ only
    \task $x = -10$ only
    \task $x = 10$ and $x = -10$
    \task $x = 50$ and $x = -50$
\end{tasks}

\begin{solution}
\textbf{Conceptual Explanation:} \\
This is a "Difference of Squares" problem, which follows the pattern $a^2 - b^2 = (a - b)(a + b)$.

\vspace{15pt}

\textbf{Calculation and Logic:} \\
1. Factor the expression: \\
   $(x - 10)(x + 10) = 0$

\vspace{10pt}

2. Solve for $x$: \\
   $x - 10 = 0 \Rightarrow x = 10$ \\
   $x + 10 = 0 \Rightarrow x = -10$

\vspace{20pt}

Answer: \colorbox{yellow!100}{C}
\end{solution}

\vspace{30pt}

% Question 19
% Category: Practice | Difficulty: Hard | Question Type: Multiple Choice
\question
Solving $x^2 + 4x - 1 = 0$ by completing the square leads to which equation?
\begin{tasks}(2)
    \task $(x + 2)^2 = 5$
    \task $(x + 2)^2 = 3$
    \task $(x + 4)^2 = 17$
    \task $(x + 2)^2 = 1$
\end{tasks}

\begin{solution}
\textbf{Conceptual Explanation:} \\
Completing the square involves creating a perfect square trinomial on one side of the equation.

\vspace{15pt}

\textbf{Calculation and Logic:} \\
1. Move the constant to the right: \\
   $x^2 + 4x = 1$

\vspace{10pt}

2. Take half of the $x$-coefficient (4), square it, and add to both sides: \\
   $(\frac{4}{2})^2 = 4$ \\
   $x^2 + 4x + 4 = 1 + 4$

\vspace{10pt}

3. Factor the left side: \\
   $(x + 2)^2 = 5$

\vspace{20pt}

Answer: \colorbox{yellow!100}{A}
\end{solution}

\vspace{30pt}

% Question 20
% Category: Practice | Difficulty: Medium | Question Type: Multiple Choice
\question
If $(x - a)(x - b) = 0$ and the solutions are $x = 7$ and $x = -2$, what is the value of $a + b$?
\begin{tasks}(4)
    \task 5
    \task 9
    \task -5
    \task 14
\end{tasks}

\begin{solution}
\textbf{Conceptual Explanation:} \\
By identifying the roots from the factored form, we can determine the values of the constants $a$ and $b$.

\vspace{15pt}

\textbf{Calculation and Logic:} \\
1. The roots are $x = a$ and $x = b$. \\
2. Therefore, $a = 7$ and $b = -2$ (or vice versa). \\
3. Sum: $7 + (-2) = 5$.

\vspace{20pt}

Answer: \colorbox{yellow!100}{A}
\end{solution}

\vspace{30pt}

% Question 21
% Category: Practice | Difficulty: Hard | Question Type: Multiple Choice
\question
Which of the following quadratic equations has no real solutions?
\begin{tasks}(2)
    \task $x^2 - 4x + 4 = 0$
    \task $x^2 + 4 = 0$
    \task $x^2 - 4 = 0$
    \task $x^2 + 4x + 3 = 0$
\end{tasks}

\begin{solution}
\textbf{Conceptual Explanation:} \\
A quadratic equation has no real solutions if the $x^2$ term plus a positive constant equals zero, because a squared number cannot be negative.

\vspace{15pt}

\textbf{Calculation and Logic:} \\
Test Choice B: $x^2 + 4 = 0$ \\
$x^2 = -4$ \\
Since no real number squared results in a negative value, this has no real solutions.

\vspace{20pt}

Answer: \colorbox{yellow!100}{B}
\end{solution}

\vspace{30pt}

% Question 22
% Category: Practice | Difficulty: Medium | Question Type: Multiple Choice
\question
What is the vertex of the parabola $y = (x - 3)^2 + 5$?
\begin{tasks}(4)
    \task $(3, 5)$
    \task $(-3, 5)$
    \task $(3, -5)$
    \task $(5, 3)$
\end{tasks}

\begin{solution}
\textbf{Conceptual Explanation:} \\
Vertex form is $y = a(x - h)^2 + k$, where $(h, k)$ is the vertex. Note that the formula uses $(x - h)$, so the sign of $h$ inside the parentheses is flipped.

\vspace{15pt}

\textbf{Calculation and Logic:} \\
1. In the equation $y = (x - 3)^2 + 5$: \\
   $h = 3$ \\
   $k = 5$

\vspace{10pt}

The vertex is $(3, 5)$.

\vspace{20pt}

Answer: \colorbox{yellow!100}{A}
\end{solution}

\vspace{30pt}

% Question 23
% Category: Practice | Difficulty: Hard | Question Type: Multiple Choice
\question
Find the roots of $x^2 - 5x = 0$.
\begin{tasks}(4)
    \task $x = 5$ only
    \task $x = 0$ and $x = -5$
    \task $x = 0$ and $x = 5$
    \task $x = 2.5$ and $x = -2.5$
\end{tasks}

\begin{solution}
\textbf{Conceptual Explanation:} \\
If an equation has a common factor in all terms, factor it out first.

\vspace{15pt}

\textbf{Calculation and Logic:} \\
1. Factor out $x$: \\
   $x(x - 5) = 0$

\vspace{10pt}

2. Set each factor to zero: \\
   $x = 0$ \\
   $x - 5 = 0 \Rightarrow x = 5$

\vspace{20pt}

Answer: \colorbox{yellow!100}{C}
\end{solution}

\vspace{30pt}

% Question 24
% Category: Practice | Difficulty: Medium | Question Type: Multiple Choice
\question
The equation $x^2 - kx + 9 = 0$ has exactly one real solution. What is a possible value for $k$?
\begin{tasks}(4)
    \task 3
    \task 6
    \task 9
    \task 0
\end{tasks}

\begin{solution}
\textbf{Conceptual Explanation:} \\
One real solution means the discriminant $b^2 - 4ac$ must equal zero.

\vspace{15pt}

\textbf{Calculation and Logic:} \\
1. Identify $a = 1, b = -k, c = 9$. \\
2. $D = (-k)^2 - 4(1)(9) = 0$ \\
3. $k^2 - 36 = 0$ \\
4. $k^2 = 36 \Rightarrow k = \pm 6$

\vspace{20pt}

Answer: \colorbox{yellow!100}{B}
\end{solution}

\vspace{30pt}

% Question 25
% Category: Practice | Difficulty: Hard | Question Type: Multiple Choice
\question
Solve for $x$: $x^2 = 3x + 10$.
\begin{tasks}(4)
    \task $x = 5, x = -2$
    \task $x = -5, x = 2$
    \task $x = 10, x = -1$
    \task $x = 5, x = 2$
\end{tasks}

\begin{solution}
\textbf{Conceptual Explanation:} \\
To solve a quadratic equation, always start by moving all terms to one side to set the equation to zero.

\vspace{15pt}

\textbf{Calculation and Logic:} \\
1. Rearrange to standard form: \\
   $x^2 - 3x - 10 = 0$

\vspace{10pt}

2. Factor: \\
   $(x - 5)(x + 2) = 0$

\vspace{10pt}

3. Solve: \\
   $x = 5, x = -2$

\vspace{20pt}

Answer: \colorbox{yellow!100}{A}
\end{solution}

% --- PART 2: Advanced Quadratic Patterns & Applications (26-50) ---

% Question 26
% Category: Practice | Difficulty: Hard | Question Type: Multiple Choice
\question
The equation $x^2 + kx + 49 = 0$ has exactly one real solution for $x$. If $k > 0$, what is the value of $k$?
\begin{tasks}(4)
    \task 7
    \task 14
    \task 28
    \task 49
\end{tasks}

\begin{solution}
\textbf{Conceptual Explanation:} \\
For a quadratic to have exactly one real solution, the discriminant must be zero ($b^2 - 4ac = 0$). This occurs when the quadratic is a perfect square trinomial.

\vspace{15pt}

\textbf{Calculation and Logic:} \\
1. Identify coefficients: $a = 1$, $b = k$, $c = 49$.

\vspace{10pt}

2. Set the discriminant to zero: \\
   $k^2 - 4(1)(49) = 0$ \\
   $k^2 - 196 = 0$

\vspace{10pt}

3. Solve for $k$: \\
   $k^2 = 196$ \\
   $k = \pm 14$

\vspace{10pt}

4. Since the problem specifies $k > 0$, the value is 14.

\vspace{20pt}

Answer: \colorbox{yellow!100}{B}
\end{solution}

\vspace{30pt}

% Question 27
% Category: Practice | Difficulty: Medium | Question Type: Multiple Choice
\question
If $(x - 10)^2 = 64$ and $x < 10$, what is the value of $x$?
\begin{tasks}(4)
    \task 2
    \task 18
    \task -2
    \task 8
\end{tasks}

\begin{solution}
\textbf{Conceptual Explanation:} \\
When solving equations involving squares, taking the square root yields two possible cases: positive and negative. The constraint $x < 10$ will help us choose the correct one.

\vspace{15pt}

\textbf{Calculation and Logic:} \\
1. Take the square root of both sides: \\
   $x - 10 = \pm 8$

\vspace{10pt}

2. Solve Case 1: \\
   $x - 10 = 8 \Rightarrow x = 18$

\vspace{10pt}

3. Solve Case 2: \\
   $x - 10 = -8 \Rightarrow x = 2$

\vspace{10pt}

4. Apply the constraint $x < 10$. The only valid solution is 2.

\vspace{20pt}

Answer: \colorbox{yellow!100}{A}
\end{solution}

\vspace{30pt}

% Question 28
% Category: Practice | Difficulty: Hard | Question Type: Multiple Choice
\question
What is the sum of the solutions to $(2x - 1)(x + 5) = 0$?
\begin{tasks}(4)
    \task 4.5
    \task -4.5
    \task -4
    \task 5.5
\end{tasks}

\begin{solution}
\textbf{Conceptual Explanation:} \\
First, find each individual solution by setting the factors to zero. Then, calculate their sum as requested.

\vspace{15pt}

\textbf{Calculation and Logic:} \\
1. Set the first factor to zero: \\
   $2x - 1 = 0 \Rightarrow 2x = 1 \Rightarrow x = 0.5$

\vspace{10pt}

2. Set the second factor to zero: \\
   $x + 5 = 0 \Rightarrow x = -5$

\vspace{10pt}

3. Calculate the sum: \\
   $0.5 + (-5) = -4.5$

\vspace{20pt}

Answer: \colorbox{yellow!100}{B}
\end{solution}

\vspace{30pt}

% Question 29
% Category: Practice | Difficulty: Hard | Question Type: Multiple Choice
\question
Which of the following is equivalent to $x^2 + 10x + 15 = 0$?
\begin{tasks}(2)
    \task $(x + 5)^2 = 10$
    \task $(x + 5)^2 = 25$
    \task $(x + 10)^2 = 85$
    \task $(x + 5)^2 = -15$
\end{tasks}

\begin{solution}
\textbf{Conceptual Explanation:} \\
To rewrite the equation, we complete the square by taking half of the $b$-coefficient, squaring it, and adding it to both sides.

\vspace{15pt}

\textbf{Calculation and Logic:} \\
1. Move the constant to the right: \\
   $x^2 + 10x = -15$

\vspace{10pt}

2. Calculate the value to complete the square: \\
   $(\frac{10}{2})^2 = 25$

\vspace{10pt}

3. Add 25 to both sides: \\
   $x^2 + 10x + 25 = -15 + 25$ \\
   $(x + 5)^2 = 10$

\vspace{20pt}

Answer: \colorbox{yellow!100}{A}
\end{solution}

\vspace{30pt}

% Question 30
% Category: Practice | Difficulty: Medium | Question Type: Multiple Choice
\question
If $x^2 - 14x + 49 = 0$, what is the value of $x$?
\begin{tasks}(4)
    \task 7
    \task -7
    \task 0
    \task 49
\end{tasks}

\begin{solution}
\textbf{Conceptual Explanation:} \\
This expression is a perfect square trinomial, following the pattern $a^2 - 2ab + b^2 = (a - b)^2$.

\vspace{15pt}

\textbf{Calculation and Logic:} \\
1. Recognize the square roots: $\sqrt{x^2} = x$ and $\sqrt{49} = 7$. \\
2. Check the middle term: $2(x)(7) = 14x$.

\vspace{10pt}

3. Factor: \\
   $(x - 7)^2 = 0$

\vspace{10pt}

4. Solve: $x = 7$.

\vspace{20pt}

Answer: \colorbox{yellow!100}{A}
\end{solution}

\vspace{30pt}

% Question 31
% Category: Practice | Difficulty: Hard | Question Type: Multiple Choice
\question
A quadratic equation is given as $ax^2 + bx + c = 0$. If $b^2 < 4ac$, which of the following is true?
\begin{tasks}(1)
    \task The equation has exactly two real solutions.
    \task The equation has exactly one real solution.
    \task The equation has no real solutions.
    \task The solutions are both positive.
\end{tasks}

\begin{solution}
\textbf{Conceptual Explanation:} \\
The number of real solutions is determined by the sign of the discriminant ($D = b^2 - 4ac$).

\vspace{15pt}

\textbf{Calculation and Logic:} \\
1. If $b^2 < 4ac$, then subtracting $4ac$ from both sides gives $b^2 - 4ac < 0$. \\
2. A negative discriminant means the quadratic formula would require taking the square root of a negative number. \\
3. This results in zero real solutions.

\vspace{20pt}

Answer: \colorbox{yellow!100}{C}
\end{solution}

\vspace{30pt}

% Question 32
% Category: Practice | Difficulty: Medium | Question Type: Multiple Choice
\question
What is the positive solution to $x^2 - \frac{1}{4} = 0$?
\begin{tasks}(4)
    \task 0.25
    \task 0.5
    \task 2
    \task 4
\end{tasks}

\begin{solution}
\textbf{Conceptual Explanation:} \\
Isolate the squared variable to find the solutions.

\vspace{15pt}

\textbf{Calculation and Logic:} \\
1. $x^2 = \frac{1}{4}$ \\
2. $x = \pm \sqrt{\frac{1}{4}}$ \\
3. $x = \pm 0.5$

\vspace{10pt}

The positive solution is 0.5.

\vspace{20pt}

Answer: \colorbox{yellow!100}{B}
\end{solution}

\vspace{30pt}

% Question 33
% Category: Practice | Difficulty: Hard | Question Type: Multiple Choice
\question
If $x = 2$ is a solution to $x^2 + kx - 10 = 0$, what is the other solution?
\begin{tasks}(4)
    \task -5
    \task 5
    \task -2
    \task 3
\end{tasks}

\begin{solution}
\textbf{Conceptual Explanation:} \\
First, find $k$ by substituting the known root. Then, factor the completed equation to find the second root.

\vspace{15pt}

\textbf{Calculation and Logic:} \\
1. Substitute $x = 2$: \\
   $(2)^2 + k(2) - 10 = 0$ \\
   $4 + 2k - 10 = 0 \Rightarrow 2k = 6 \Rightarrow k = 3$.

\vspace{10pt}

2. Rewrite the equation: \\
   $x^2 + 3x - 10 = 0$

\vspace{10pt}

3. Factor: \\
   $(x + 5)(x - 2) = 0$

\vspace{10pt}

4. Identify the roots: $x = -5$ and $x = 2$.

\vspace{20pt}

Answer: \colorbox{yellow!100}{A}
\end{solution}

\vspace{30pt}

% Question 34
% Category: Practice | Difficulty: Medium | Question Type: Multiple Choice
\question
Which of the following expressions is a factor of $x^2 + 5x - 6$?
\begin{tasks}(4)
    \task $x + 2$
    \task $x - 3$
    \task $x + 6$
    \task $x - 6$
\end{tasks}

\begin{solution}
\textbf{Conceptual Explanation:} \\
We need two numbers that multiply to $-6$ and add up to $5$.

\vspace{15pt}

\textbf{Calculation and Logic:} \\
1. Multipliers of $-6$: $(6, -1)$, $(-6, 1)$, $(3, -2)$, $(-3, 2)$. \\
2. Sum of $(6, -1) = 5$.

\vspace{10pt}

3. Factored form: $(x + 6)(x - 1)$. \\
   $x + 6$ is the factor listed in the choices.

\vspace{20pt}

Answer: \colorbox{yellow!100}{C}
\end{solution}

\vspace{30pt}

% Question 35
% Category: Practice | Difficulty: Hard | Question Type: Multiple Choice
\question
If $2(x - 3)^2 = 18$, what is the sum of the possible values of $x$?
\begin{tasks}(4)
    \task 0
    \task 3
    \task 6
    \task 12
\end{tasks}

\begin{solution}
\textbf{Conceptual Explanation:} \\
Isolate the squared term before taking the square root.

\vspace{15pt}

\textbf{Calculation and Logic:} \\
1. Divide by 2: $(x - 3)^2 = 9$. \\
2. Square root: $x - 3 = \pm 3$.

\vspace{10pt}

3. Values of $x$: \\
   $x = 3 + 3 = 6$ \\
   $x = 3 - 3 = 0$

\vspace{10pt}

4. Sum: $6 + 0 = 6$.

\vspace{20pt}

Answer: \colorbox{yellow!100}{C}
\end{solution}

\vspace{30pt}

% Question 36
% Category: Practice | Difficulty: Hard | Question Type: Multiple Choice
\question
For what value of $c$ will the equation $x^2 - 12x + c = 0$ have only one solution?
\begin{tasks}(4)
    \task 6
    \task 12
    \task 36
    \task 144
\end{tasks}

\begin{solution}
\textbf{Conceptual Explanation:} \\
In a perfect square trinomial $x^2 + bx + c$, the value of $c$ is $(\frac{b}{2})^2$.

\vspace{15pt}

\textbf{Calculation and Logic:} \\
1. $b = -12$. \\
2. $c = (\frac{-12}{2})^2 = (-6)^2 = 36$.

\vspace{20pt}

Answer: \colorbox{yellow!100}{C}
\end{solution}

\vspace{30pt}

% Question 37
% Category: Practice | Difficulty: Medium | Question Type: Multiple Choice
\question
What are the solutions to $x^2 + 8x + 12 = 0$?
\begin{tasks}(4)
    \task $x = 6, x = 2$
    \task $x = -6, x = -2$
    \task $x = -4, x = -4$
    \task $x = -12, x = -1$
\end{tasks}

\begin{solution}
\textbf{Conceptual Explanation:} \\
Find factors of 12 that sum to 8.

\vspace{15pt}

\textbf{Calculation and Logic:} \\
1. Factors of 12: $(6, 2)$. \\
2. Sum: $6 + 2 = 8$.

\vspace{10pt}

3. Factored form: $(x + 6)(x + 2) = 0$. \\
4. Roots: $x = -6$ and $x = -2$.

\vspace{20pt}

Answer: \colorbox{yellow!100}{B}
\end{solution}

\vspace{30pt}

% Question 38
% Category: Practice | Difficulty: Hard | Question Type: Multiple Choice
\question
If $x^2 - y^2 = 24$ and $x + y = 6$, what is the value of $x - y$?
\begin{tasks}(4)
    \task 4
    \task 18
    \task 144
    \task 2
\end{tasks}

\begin{solution}
\textbf{Conceptual Explanation:} \\
Use the difference of squares identity: $x^2 - y^2 = (x - y)(x + y)$.

\vspace{15pt}

\textbf{Calculation and Logic:} \\
1. Substitute the known values: \\
   $24 = (x - y)(6)$

\vspace{10pt}

2. Solve for $(x - y)$: \\
   $x - y = \frac{24}{6} = 4$.

\vspace{20pt}

Answer: \colorbox{yellow!100}{A}
\end{solution}

\vspace{30pt}

% Question 39
% Category: Practice | Difficulty: Medium | Question Type: Multiple Choice
\question
Which of the following equations has roots $x = 3$ and $x = -5$?
\begin{tasks}(2)
    \task $(x + 3)(x - 5) = 0$
    \task $(x - 3)(x + 5) = 0$
    \task $(x - 3)(x - 5) = 0$
    \task $(x + 3)(x + 5) = 0$
\end{tasks}

\begin{solution}
\textbf{Conceptual Explanation:} \\
To form an equation from roots $r_1$ and $r_2$, use the factors $(x - r_1)$ and $(x - r_2)$.

\vspace{15pt}

\textbf{Calculation and Logic:} \\
1. Root $x = 3 \Rightarrow$ Factor $(x - 3)$. \\
2. Root $x = -5 \Rightarrow$ Factor $(x - (-5)) = (x + 5)$.

\vspace{10pt}

3. Equation: $(x - 3)(x + 5) = 0$.

\vspace{20pt}

Answer: \colorbox{yellow!100}{B}
\end{solution}

\vspace{30pt}

% Question 40
% Category: Practice | Difficulty: Hard | Question Type: Multiple Choice
\question
What is the value of the discriminant for $3x^2 - 2x + 1 = 0$?
\begin{tasks}(4)
    \task -8
    \task 16
    \task 4
    \task 0
\end{tasks}

\begin{solution}
\textbf{Conceptual Explanation:} \\
The discriminant is $b^2 - 4ac$.

\vspace{15pt}

\textbf{Calculation and Logic:} \\
1. $a = 3, b = -2, c = 1$. \\
2. $D = (-2)^2 - 4(3)(1)$ \\
3. $D = 4 - 12 = -8$.

\vspace{20pt}

Answer: \colorbox{yellow!100}{A}
\end{solution}

\vspace{30pt}

% Question 41
% Category: Practice | Difficulty: Hard | Question Type: Multiple Choice
\question
If $2x^2 + 8x = 0$, what are the possible values for $x$?
\begin{tasks}(4)
    \task 0 and 4
    \task 0 and -4
    \task 4 and -4
    \task 2 and -8
\end{tasks}

\begin{solution}
\textbf{Conceptual Explanation:} \\
Factor out the greatest common factor (GCF).

\vspace{15pt}

\textbf{Calculation and Logic:} \\
1. GCF is $2x$. \\
2. $2x(x + 4) = 0$.

\vspace{10pt}

3. Set factors to zero: \\
   $2x = 0 \Rightarrow x = 0$ \\
   $x + 4 = 0 \Rightarrow x = -4$

\vspace{20pt}

Answer: \colorbox{yellow!100}{B}
\end{solution}

\vspace{30pt}

% Question 42
% Category: Practice | Difficulty: Medium | Question Type: Multiple Choice
\question
What is the vertex form of $y = x^2 - 6x + 13$?
\begin{tasks}(2)
    \task $y = (x - 3)^2 + 4$
    \task $y = (x - 3)^2 + 13$
    \task $y = (x + 3)^2 + 4$
    \task $y = (x - 6)^2 + 7$
\end{tasks}

\begin{solution}
\textbf{Conceptual Explanation:} \\
Complete the square to transition from standard form to vertex form.

\vspace{15pt}

\textbf{Calculation and Logic:} \\
1. $x^2 - 6x + 9 - 9 + 13$ \\
2. $(x - 3)^2 + 4$

\vspace{20pt}

Answer: \colorbox{yellow!100}{A}
\end{solution}

\vspace{30pt}

% Question 43
% Category: Practice | Difficulty: Medium | Question Type: Multiple Choice
\question
Solve for $x$: $(x + 2)^2 = 0$.
\begin{tasks}(4)
    \task 2
    \task -2
    \task 0
    \task $\pm 2$
\end{tasks}

\begin{solution}
\textbf{Conceptual Explanation:} \\
A squared expression equals zero only when the expression itself is zero.

\vspace{15pt}

\textbf{Calculation and Logic:} \\
1. $x + 2 = 0$ \\
2. $x = -2$.

\vspace{20pt}

Answer: \colorbox{yellow!100}{B}
\end{solution}

\vspace{30pt}

% Question 44
% Category: Practice | Difficulty: Hard | Question Type: Multiple Choice
\question

Which equation matches the graph shown above?
\begin{tasks}(2)
    \task $y = (x - 1)(x - 3)$
    \task $y = (x + 1)(x + 3)$
    \task $y = (x - 2)^2 + 1$
    \task $y = x^2 - 4$
\end{tasks}

\begin{solution}
\textbf{Conceptual Explanation:} \\
Identify the roots (where the graph touches the $x$-axis) and use them to find the factored form.

\vspace{15pt}

\textbf{Calculation and Logic:} \\
1. Roots are $x = 1$ and $x = 3$. \\
2. Factored form: $(x - 1)(x - 3)$.

\vspace{20pt}

Answer: \colorbox{yellow!100}{A}
\end{solution}

\vspace{30pt}

% Question 45
% Category: Practice | Difficulty: Hard | Question Type: Multiple Choice
\question
What is the value of $x$ for the equation $x^2 + 1 = 0$?
\begin{tasks}(4)
    \task 1
    \task -1
    \task No real solution
    \task 0
\end{tasks}

\begin{solution}
\textbf{Conceptual Explanation:} \\
Check the discriminant or solve directly to see if real roots exist.

\vspace{15pt}

\textbf{Calculation and Logic:} \\
1. $x^2 = -1$ \\
2. Since no real number squared can be negative, there is no real solution.

\vspace{20pt}

Answer: \colorbox{yellow!100}{C}
\end{solution}

\vspace{30pt}

% Question 46
% Category: Practice | Difficulty: Medium | Question Type: Multiple Choice
\question
If $f(x) = x^2 - 4x + 4$, what is $f(5)$?
\begin{tasks}(4)
    \task 9
    \task 1
    \task 5
    \task 25
\end{tasks}

\begin{solution}
\textbf{Conceptual Explanation:} \\
Substitute the given value of $x$ into the function.

\vspace{15pt}

\textbf{Calculation and Logic:} \\
1. $5^2 - 4(5) + 4$ \\
2. $25 - 20 + 4 = 9$.

\vspace{20pt}

Answer: \colorbox{yellow!100}{A}
\end{solution}

\vspace{30pt}

% Question 47
% Category: Practice | Difficulty: Hard | Question Type: Multiple Choice
\question
Which of the following is the quadratic formula?
\begin{tasks}(2)
    \task $x = \frac{-b \pm \sqrt{b^2 - 4ac}}{2a}$
    \task $x = \frac{b \pm \sqrt{b^2 - 4ac}}{2a}$
    \task $x = \frac{-b \pm \sqrt{b^2 + 4ac}}{2a}$
    \task $x = -b \pm \sqrt{b^2 - 4ac}$
\end{tasks}

\begin{solution}
\textbf{Conceptual Explanation:} \\
The quadratic formula is a universal method to find the roots of any quadratic equation.

\vspace{15pt}

\textbf{Calculation and Logic:} \\
The standard formula is $x = \frac{-b \pm \sqrt{b^2 - 4ac}}{2a}$.

\vspace{20pt}

Answer: \colorbox{yellow!100}{A}
\end{solution}

\vspace{30pt}

% Question 48
% Category: Practice | Difficulty: Medium | Question Type: Multiple Choice
\question
What is the $y$-intercept of the function $y = x^2 - 5x + 6$?
\begin{tasks}(4)
    \task (0, 6)
    \task (6, 0)
    \task (0, -5)
    \task (2, 3)
\end{tasks}

\begin{solution}
\textbf{Conceptual Explanation:} \\
The $y$-intercept occurs where $x = 0$.

\vspace{15pt}

\textbf{Calculation and Logic:} \\
1. Substitute $x = 0$: $y = 0^2 - 5(0) + 6$. \\
2. $y = 6$.

\vspace{20pt}

Answer: \colorbox{yellow!100}{A}
\end{solution}

\vspace{30pt}

% Question 49
% Category: Practice | Difficulty: Hard | Question Type: Multiple Choice
\question
Solve for $x$: $x^2 + 6x = 7$.
\begin{tasks}(4)
    \task $x = 1, x = -7$
    \task $x = -1, x = 7$
    \task $x = 1, x = 7$
    \task $x = -1, x = -7$
\end{tasks}

\begin{solution}
\textbf{Conceptual Explanation:} \\
Set the equation to zero before factoring.

\vspace{15pt}

\textbf{Calculation and Logic:} \\
1. $x^2 + 6x - 7 = 0$ \\
2. $(x + 7)(x - 1) = 0$ \\
3. $x = -7, x = 1$.

\vspace{20pt}

Answer: \colorbox{yellow!100}{A}
\end{solution}

\vspace{30pt}

% Question 50
% Category: Practice | Difficulty: Hard | Question Type: Multiple Choice
\question
The vertex of a parabola is $(h, k)$. If the equation is $y = (x + 2)^2 - 9$, what are the $x$-intercepts?
\begin{tasks}(4)
    \task 1 and -5
    \task -1 and 5
    \task 3 and -3
    \task 2 and -9
\end{tasks}

\begin{solution}
\textbf{Conceptual Explanation:} \\
$x$-intercepts occur when $y = 0$.

\vspace{15pt}

\textbf{Calculation and Logic:} \\
1. $0 = (x + 2)^2 - 9$ \\
2. $9 = (x + 2)^2$ \\
3. $\pm 3 = x + 2$

\vspace{10pt}

4. Case 1: $x + 2 = 3 \Rightarrow x = 1$. \\
5. Case 2: $x + 2 = -3 \Rightarrow x = -5$.

\vspace{20pt}

Answer: \colorbox{yellow!100}{A}
\end{solution}
% =============================================================
% Question End
% =============================================================

\end{questions}
\begin{center}
\rule{.5\textwidth}{1pt}
\end{center}
\end{document}
